\documentclass[../main.tex]{subfiles}
\begin{document}

%========================================================================================
% FILENAME: 07_discipline.tex (v1.0.0)
%
% §7 Market Discipline.  §7.1 three feedback forces + cost-convergence table; §7.2
% fee-channel anti-correlation (uses the surface handle \confprice, always applied).
%
% Proof budget: one \proofsketch margin (anti-correlation).  No \interfaceref.  Silences
% applied; cycle only; "actor" -> "agent".
%
% Change history lives in version control (see the versioning policy in main.tex).
%========================================================================================

\section{Market Discipline}
\label{sec:attestation:discipline}

Once external holders exist, three forces create negative feedback around burning; none
exists in the single-holder regime. This is the market property that activates the
\emph{net} costliness the protocol alone cannot supply
(\zcref{intuit:costliness:design-problem}).

% ═══════════════════════════════════════════════════════════════════════
\subsection{Three Feedback Forces}
\label{sec:discipline:forces}
% ═══════════════════════════════════════════════════════════════════════

\textbf{Redemption drainage.} Burns raise $\ratefloor{}$; other holders redeem at the new
floor, extracting reserve units from $\reservepool{}$ and reducing future burn
productivity. \textbf{Acquisition cost.} Any agent who wants to burn must deposit reserve
units or acquire \liveclass{} units on a secondary market at the market price.
\textbf{Speculative holding.} Holders expecting future burns to raise $\ratefloor{}$ delay
redemption, tightening liquid supply and raising acquisition cost.

\begin{equation*}
    \text{Burn}
    \xrightarrow{\text{raises } \ratefloor{}}
    \text{Redemption}
    \xrightarrow{\text{drains } \reservepool{}}
    \text{Reduced burn productivity}
\end{equation*}
\begin{equation*}
    \text{Burn}
    \xrightarrow{\text{raises } \ratefloor{}}
    \text{Speculative holding}
    \xrightarrow{\text{tightens supply}}
    \text{Higher acquisition cost}
\end{equation*}

Both loops are negative: burning triggers responses that make further burning less
productive or more expensive. Neither exists in the single-holder regime.

\begin{table}[h!]
    \centering
    \caption{Cost convergence under external deposit dilution. Independent of $\livesplit$
    (\zcref{thm:costliness:zeta-independence}).}
    \label{tab:discipline:cost-convergence}
    \begin{tabular}{@{}c c c l@{}}
        \toprule
        $D/\genesisreserve$ & $\supplyfrac{}$ & Net cost / unit & \textbf{Regime} \\
        \midrule
        $0$         & $1.00$    & $0$    & Monopoly \\
        $1.0$       & $0.50$    & $0.50$ & Transitional \\
        $10.0$      & $0.09$    & $0.91$ & Market discipline \\
        $\to\infty$ & $\to 0$   & $\to 1$ & Full costliness \\
        \bottomrule
    \end{tabular}
\end{table}

% ═══════════════════════════════════════════════════════════════════════
\subsection{Fee-Channel Anti-Correlation}
\label{sec:discipline:anti-correlation}
% ═══════════════════════════════════════════════════════════════════════

\begin{proposition}[Fee--Time-Locked-Class Channel Anti-Correlation]
    \label{prop:discipline:anti-correlation}
    \proofsketch{$\confprice{d}$ rises and $\supplyfrac{}(d)$ falls in $d$; the product is
    bounded.}%
    The operator receives $\feerate(1-\livesplit)\Delta\totsupply{}$ fee \tlockclass{} units
    per cycle, which are transferable. Selling them for \liveclass{} units and burning the
    acquired units would circumvent the $\livesplit$-suppression of
    \zcref{prop:containment:seigniorage-suppression}. The channel is self-resolving: with
    $d = D/\genesisreserve$, the time-locked-class secondary-market confidence price
    $\confprice{d}$ approaches live-class parity as $d \to \infty$, while the genesis supply
    dominance $\supplyfrac{}(d) = 1/(1+d) \to 0$. Since $\confprice{d}$ is monotone
    increasing and $\supplyfrac{}(d)$ monotone decreasing in $d$: when the channel is
    operative ($\confprice{d}$ near parity, $d$ large), $\supplyfrac{} \approx 0$ and burns
    are genuinely costly; when the genesis advantage exists ($\supplyfrac{} \approx 1$, $d$
    small), $\confprice{d}$ is deeply discounted and the channel yields negligible
    \liveclass{} units per \tlockclass{} unit sold.
\end{proposition}

\begin{remark}[Transient Decorrelation]
    \label{rem:discipline:transient-decorrelation}
    The anti-correlation is a steady-state relationship; during transient liquidity events
    it may temporarily break. The containment bounds (\zcref{sec:attestation:containment}) bind from
    genesis regardless of secondary-market conditions.
\end{remark}

%========================================================================================
% END OF FILE: 07_discipline.tex (v1.0.0)
%========================================================================================

\end{document}
